% Generated draft: requires mathematical and source review.
\documentclass[11pt]{article}
\usepackage[margin=25mm]{geometry}
\usepackage{fontspec}
\setmainfont{Latin Modern Roman}
\usepackage{amsmath,amssymb,mathtools,unicode-math}
\setmathfont{Latin Modern Math}
\usepackage{enumitem,needspace,xurl,hyperref}
\setlength{\parindent}{0pt}
\setlength{\parskip}{5pt}
\setlength{\emergencystretch}{3em}
\newcommand{\N}{\mathbb{N}}
\newcommand{\Z}{\mathbb{Z}}
\newcommand{\Q}{\mathbb{Q}}
\newcommand{\R}{\mathbb{R}}
\newcommand{\C}{\mathbb{C}}
\newcommand{\F}{\mathbb{F}}
\newcommand{\A}{\mathbb{A}}
\newcommand{\PP}{\mathbb{P}}
\newcommand{\E}{\mathbb{E}}
\newcommand{\eps}{\varepsilon}
\newcommand{\dd}{\,\mathrm{d}}
\begin{document}
\section*{UnsolvedMath ID: 30005951}
\textbf{Source problem code:} catalogue entry
\subsection{Definitions and mathematical statement}
Write \(R=\F[U]\), where the coefficient field \(\F\) is not specified in the supplied record. Let \(M\) be an \(R\)-module and let \(M_{\mathrm{red}}\) denote its reduced module in the terminology of the source. The source does not define this reduced module or specify its relationship to the torsion submodule. Assuming that some power of \(U\) annihilates \(M_{\mathrm{red}}\), put
\[
\ell=\min\{k\in\Z_{\geq 0}:U^kM_{\mathrm{red}}=0\}.
\]
The supplied definition says that \(M\) satisfies the \emph{strong geography restriction} if there exists an \(R\)-module \(N\) such that
\[
M\cong N\oplus\bigoplus_{j=1}^{\ell}R/(U^j).
\]
This expresses the displayed direct-summand condition at the level of \(R\)-modules; the record specifies no grading requirement or grading shifts. When \(\ell=0\), the sum is interpreted as the zero module, so the condition is vacuous. If no power of \(U\) annihilates \(M_{\mathrm{red}}\), the displayed minimum does not exist and this definition does not apply without further conventions.

A rational homology three-sphere is a closed connected three-manifold \(Y\) with \(H_*(Y;\Q)\cong H_*(S^3;\Q)\). The notation \(HF^-(Y)\) denotes its minus Heegaard Floer homology module; coefficient, orientation, grading, and \(\operatorname{Spin}^c\)-summation conventions are not specified in this record.

The extraction reports the following as an already proved theorem: if a rational homology three-sphere \(Y\) is obtained either by surgery on a knot in \(S^3\), or by large surgery on a link in \(S^3\), then \(HF^-(Y)\) satisfies the strong geography restriction above. The surgery coefficients and the meaning and threshold of ``large'' are absent, as is the reference completing the phrase ``was proved in.'' Consequently, this is a qualified transcription of a source-reported theorem, not a fully recovered theorem with all hypotheses fixed.

The surrounding discussion reports that the L-space conjecture holds for graph manifolds, some families of knot surgeries, and some branched coverings. It presents a logical alternative: either that conjecture fails, or not every module can occur as the torsion module of the Heegaard Floer homology of a rational homology sphere. It also describes the strong restriction as potentially stronger than a restriction attributed to Lin's theorem in conjunction with the L-space conjecture. Neither the relevant module class nor Lin's theorem is stated here, so this discussion cannot be converted into a precise additional quantified assertion.

No separate open problem is recoverable from this record. Its supplied verification marks the statement as unrecoverable and says that the subsequent strong-geography question is retained completely in source record \texttt{OWR-14298580-008}, whose contents are not supplied. In particular, asserting the strong geography restriction for every rational homology three-sphere would introduce a target not explicitly stated in this record.
\subsection{Short English statement}
The source defines a strong geography restriction by requiring a direct summand containing one \(\F[U]/(U^j)\) for each \(1\leq j\leq\ell\), where \(\ell\) is the nilpotence length of the reduced module. It reports this restriction as proved for specified surgery classes, but omits essential conventions and surgery hypotheses. The registry marks the separate open question as unrecoverable here and points to \texttt{OWR-14298580-008}. This is a statement-only draft requiring human mathematical review.
\subsection{Sources}
\begin{itemize}
\item Supplied local UnsolvedMath registry record, permanent ID \(30005951\), exact source problem code \texttt{catalogue entry}, title ``Strong Geography Restrictions in Heegaard Floer Homology,'' in the Oberwolfach Reports Open Problems collection. The registry supplies status ``open'' but statement status ``unrecoverable''; these are supplied curation labels, not independently verified mathematical conclusions. No proposer is identified; the proposed year is recorded as 2024.
\item Supplied report citation: \emph{Topologie} (2025), Oberwolfach Reports, pp.~1913--1974, DOI \texttt{10.4171/owr/2024/34}. This is also the literature source labelled ``original OWR report'': \url{https://doi.org/10.4171/owr/2024/34}. The record attributes its material to workshop reports and open-problem discussions; the speaker and the theorem's precise internal reference are not supplied.
\item Supplied literature source labelled ``Alfieri--Binns primary preprint'': \url{https://arxiv.org/abs/2404.00490}. No title, theorem number, or explicit correspondence between this preprint and the extracted theorem is supplied.
\item Literature assessment dated exactly \texttt{2026-08-22}: the supplied curation labels the status open and states that neither the universal L-space conjecture nor the universal geography restriction is settled. This is supplied dated curation derived from the OWR material, not an independently verified assessment. The record itself requests independent verification of time-sensitive status and cited claims.
\item The supplied statement-verification note directs the reader to \texttt{OWR-14298580-008} for the complete strong-geography question. That separate record has not been supplied here.
\end{itemize}
\end{document}
